\section{Weekly Summary}

\begin{tcolorbox}[colback=yellow!10,colframe=black!50,title=AI-Generated Content]
\noindent The summary below was written by Claude (Anthropic) from that week's regression outputs and series descriptions. It has not been reviewed or fact-checked by a human, and should be read as an automated, informal recap -- not analysis.
\end{tcolorbox}

Welcome back to the weekly ritual where we throw two unrelated economic time series into a blender and solemnly report whatever sludge comes out. This week the random number generator, in its infinite wisdom, decided to introduce goods-producing workers in Springfield, Massachusetts to a nursing home price index, bank returns in Missouri to IT orders, and residential electric power prices to America's spending on burgers and beer. A truly curated salon of strangers who have never met.

Monday kicked things off with a bang: Springfield's goods-producing employment versus proprietary and government nursing home prices posted a raw R-squared of 0.78, which, if you squint and ignore everything you know, looks downright impressive. Then we detrended it, and the relationship evaporated so completely that the R-squared slumped to 0.02 and the p-value climbed to a majestic 0.999 -- as close to "absolutely nothing" as statistics will politely allow. In other words, factory jobs in Springfield were quietly shrinking while nursing home prices were quietly rising, and our regression mistook two people walking in opposite directions for a passionate romance. Classic trend mirage, and exactly the kind of thing this project exists to point and laugh at.

Wednesday gave us IT manufacturing orders versus the return on equity for Missouri banks, a series so beloved it was literally discontinued. The raw version was significant, which is unsurprising, but here's the mildly interesting part: the relationship actually survived detrending, holding on with a p-value around 0.01. It's weak -- an R-squared of 0.06 is not going to get anyone tenure -- but it technically clung to life after we stripped out the shared drift, which puts it in the rarer "huh, maybe there's a faint something there" category rather than the usual "pure coincidence" bin. I wouldn't bet my retirement on Missouri banking returns forecasting the nation's IT appetite, but I'll grudgingly tip my hat.

Then Thursday served up the week's grand finale: residential electric power prices in the East North Central census division versus national spending on food services and drinking places. The raw fit was a glorious 0.74, and even after detrending it held up at an R-squared of 0.21 with a p-value small enough to make a p-value enthusiast weep. Does your bar tab rise with your electric bill? Possibly it's inflation gently inflating everything at once, or possibly electricity and eating out are both just things people do more of as life gets pricier -- but this one at least refused to die when we took away its trend crutch. As always, the mandatory disclaimer: these are randomly paired series, nobody is claiming anyone causes anyone, and spurious correlation remains the house specialty. If anything surprised us this week, it's that two of our three blind dates actually exchanged phone numbers.
